% ---------------------------------------------------------------------------
% Section « Déploiement » (phase 2 MLOps) — incluse depuis reports/report.tex
% via % ---------------------------------------------------------------------------
% Section « Déploiement » (phase 2 MLOps) — incluse depuis reports/report.tex
% via % ---------------------------------------------------------------------------
% Section « Déploiement » (phase 2 MLOps) — incluse depuis reports/report.tex
% via % ---------------------------------------------------------------------------
% Section « Déploiement » (phase 2 MLOps) — incluse depuis reports/report.tex
% via \include{sections/deploiement_aws}, entre sections/mlops et
% sections/conclusion.
% Utilise uniquement des paquets déjà chargés par report.tex :
% tikz (positioning, arrows.meta), booktabs, tabularx, xcolor (couleurs liora*).
% ---------------------------------------------------------------------------

\section{Déploiement sur AWS}
\label{sec:deploiement-aws}

Le service de classification tourne sur une machine virtuelle AWS partagée par
l'équipe, en région Paris (\texttt{eu-west-3}). Toute la chaîne MLOps y est
conteneurisée : suivi des expériences et registre de modèles (MLflow), stockage
des artefacts (MinIO), API de prédiction authentifiée (BentoML), supervision
(Prometheus, Grafana). Les jeux de données sont versionnés avec DVC dans un bucket
S3, et le code est contrôlé puis empaqueté par GitHub Actions. La machine n'expose
aucun port sur Internet : on n'y accède qu'à travers AWS Systems Manager, avec les
droits IAM de chacun.

\subsection{Architecture}

\begin{figure}[H]
\centering
\resizebox{\linewidth}{!}{%
\begin{tikzpicture}[
  font=\sffamily\scriptsize,
  >={Stealth[length=2.2mm]},
  box/.style={draw=lioraink, rounded corners=2pt, fill=white, align=center,
              line width=0.6pt, inner sep=3pt},
  svc/.style={box, minimum width=2.8cm, minimum height=1.5cm, text width=2.6cm},
  ctr/.style={box, minimum width=4.2cm, minimum height=2.0cm, text width=4.0cm},
  ext/.style={box, fill=lioraline!45, minimum height=1.3cm},
  flow/.style={->, draw=lioraink, line width=0.6pt},
  sec/.style={->, draw=liorared, line width=0.8pt},
  lbl/.style={font=\sffamily\tiny, text=lioragrey, align=center, fill=white, inner sep=1pt},
]
% --- Acteurs externes --------------------------------------------------------
\node[ext, minimum width=5.9cm, text width=5.7cm] (team) at (3.35,12.65)
  {\textbf{Postes de l'équipe} (3 membres)\\AWS CLI \texttt{aws login} + MFA,
   plugin Session Manager};
\node[ext, minimum width=5.9cm, text width=5.7cm] (gh) at (9.55,12.65)
  {\textbf{GitHub} \texttt{jul26\_bds\_blood\_cells}\\Actions : CI (ruff, pytest),
   image de l'API testée puis publiée};

% --- Compte AWS -------------------------------------------------------------
\draw[lioraink, line width=0.8pt, rounded corners=4pt] (0,0) rectangle (16,11.2);
\node[anchor=north east, font=\sffamily\small\bfseries, text=lioraink] at (15.85,11.12)
  {AWS \textperiodcentered{} région eu-west-3 (Paris)};

\node[svc] (iam)   at (1.8,9.6)  {\textbf{IAM}\\groupe \texttt{blood-cells-team}\\garde-fous :\\région, suppression};
\node[svc] (ssm)   at (4.9,9.6)  {\textbf{Systems Manager}\\Session Manager\\(shell, tunnels)};
\node[svc] (ecr)   at (8.0,9.6)  {\textbf{ECR}\\registre d'images\\(image de l'API)};
\node[svc] (s3)    at (11.1,9.6) {\textbf{S3}\\jeux de données\\versionnés (DVC)};
\node[svc] (sched) at (14.2,9.6) {\textbf{EventBridge Scheduler}\\arrêt de la VM à 22 h};

% --- Machine virtuelle --------------------------------------------------------
\draw[liorared, line width=0.9pt, rounded corners=4pt] (0.3,0.3) rectangle (15.7,8.0);
\node[anchor=north west, align=left, text=lioraink] at (0.45,7.95)
  {\textbf{EC2 \texttt{blood-cells-vm}} \textperiodcentered{} m7i-flex.large (2 vCPU, 8 Go, 80 Go)
   \textperiodcentered{} Ubuntu 24.04 \textperiodcentered{} \textcolor{liorared}{aucun port entrant}
   \textperiodcentered{} rôle \texttt{blood-cells-ec2-role}};

\draw[lioragrey, dashed, rounded corners=3pt] (0.6,0.6) rectangle (15.4,7.1);
\node[anchor=north west, text=lioragrey] at (0.7,7.05) {Docker Compose};

\node[ctr] (api)   at (3.1,5.2)  {\textbf{API BentoML} (8000)\\\texttt{/login} (JWT),
  \texttt{/predict}, \texttt{/models}, \texttt{/metrics}};
\node[ctr] (mlf)   at (7.9,5.2)  {\textbf{MLflow} (5001)\\suivi des expériences,\\registre, alias \emph{champion}};
\node[ctr] (minio) at (12.7,5.2) {\textbf{MinIO} (9000, 9001)\\stockage S3 des artefacts MLflow};
\node[ctr] (prom)  at (3.1,2.45) {\textbf{Prometheus} (9090)\\collecte des métriques,\\4 règles d'alerte};
\node[ctr] (train) at (7.9,2.45) {\textbf{Entraînement}\\\texttt{src.experiments.run}\\lit \texttt{data/} (PBC, LMU)};
\node[ctr] (graf)  at (12.7,2.45) {\textbf{Grafana} (9091)\\tableau de bord de l'API};

% --- Flux -------------------------------------------------------------------
\draw[sec]  (team.south -| iam)  -- node[lbl, pos=0.45] {authentification} (iam.north);
\draw[sec]  (team.south -| ssm)  -- node[lbl, pos=0.45] {session chiffrée} (ssm.north);
\draw[sec]  (gh.south -| ecr)    -- node[lbl, pos=0.45] {OIDC, \texttt{docker push}} (ecr.north);

\draw[sec]  (ssm.south)   -- node[lbl] {shell et\\tunnels de ports} (ssm.south |- 0,8.0);
\draw[flow] (ecr.south)   -- node[lbl] {\texttt{docker pull}} (ecr.south |- 0,8.0);
\draw[flow] (s3.south)    -- node[lbl] {\texttt{dvc pull}} (s3.south |- 0,8.0);
\draw[flow] (sched.south) -- node[lbl] {arrêt} (sched.south |- 0,8.0);

\draw[flow] (api)   -- node[lbl, above=1pt] {charge le\\\emph{champion}} (mlf);
\draw[flow] (mlf)   -- node[lbl, above=1pt] {artefacts} (minio);
\draw[flow] (train) -- node[lbl] {runs, modèles} (mlf);
\draw[flow] (prom)  -- node[lbl] {lit \texttt{/metrics}} (api);
\draw[flow] (graf.south) -- ++(0,-0.45) -| node[lbl, pos=0.25] {requêtes} (prom.south);
\end{tikzpicture}}
\caption{Architecture déployée. En rouge, les accès contrôlés par IAM : personne
n'ouvre de port réseau, tout passe par AWS Systems Manager ou par la fédération
OIDC de GitHub. En pointillés, les services de la pile Docker Compose.}
\label{fig:architecture-aws}
\end{figure}

\subsection{Services et accès}

\begin{table}[H]
\centering
\small
\begin{tabularx}{\linewidth}{@{}l l X@{}}
\toprule
\textbf{Service} & \textbf{Port (VM)} & \textbf{Rôle} \\
\midrule
API BentoML & 8000 & prédiction ; \texttt{/login} délivre un jeton JWT valable 1 h, exigé par \texttt{/predict} \\
MLflow      & 5001 & suivi des entraînements et registre ; l'API sert la version d'alias \emph{champion} \\
MinIO       & 9000, 9001 & stockage objet des artefacts MLflow (API S3, console) \\
Prometheus  & 9090 & collecte toutes les 15 s les métriques exposées par l'API \\
Grafana     & 9091 & tableaux de bord de supervision \\
\bottomrule
\end{tabularx}
\caption{Services de la pile Docker Compose. Les ports ne sont joignables que
depuis la VM ou par un tunnel Session Manager ouvert depuis le poste d'un membre.}
\label{tab:services-aws}
\end{table}

\subsection{Sécurité}

\begin{itemize}[leftmargin=*]
  \item \textbf{Identités.} Chaque membre a son utilisateur IAM, membre du groupe
        \texttt{blood-cells-team}, avec MFA. La CLI se connecte par
        \texttt{aws login} : aucune clé d'accès permanente n'est stockée sur les postes.
  \item \textbf{Garde-fous.} Une politique de groupe refuse toute action hors
        des régions Paris et \texttt{us-east-1} (services globaux), ainsi que la
        suppression de la VM, de son disque et du bucket (étiquette \texttt{Protected}).
  \item \textbf{Réseau.} Le groupe de sécurité de la VM n'a aucune règle entrante ;
        ni SSH ni clé SSH. Shell et tunnels passent par Session Manager.
  \item \textbf{Secrets.} Mots de passe MinIO et API, clé de signature JWT :
        générés aléatoirement dans un fichier \texttt{.env} présent uniquement sur
        la VM (droits 600, ignoré par Git). La VM lit le dépôt avec un jeton
        GitHub en lecture seule. GitHub Actions publie dans ECR par fédération
        OIDC : aucun secret AWS n'est stocké dans GitHub, et seule la branche
        \texttt{main} peut endosser le rôle de publication.
  \item \textbf{Coût.} La VM est arrêtée automatiquement chaque soir à 22 h, et des
        alertes de facturation sont actives.
\end{itemize}

\subsection{Données, intégration continue et supervision}

\paragraph{Données versionnées.} Les deux jeux d'images (PBC et LMU) sont suivis par
DVC : le dépôt Git ne contient que deux petits fichiers \texttt{.dvc}, qui décrivent
le contenu exact de chaque dossier, tandis que les images sont stockées dans
\texttt{s3://blood-cells-123456789012/dvc} (35\,451 objets, 12,0~Go, les fichiers
identiques n'étant stockés qu'une fois). Chaque commit référence donc la version
précise des données ; sur la VM, \texttt{dvc pull} les restaure en quelques secondes
grâce au rôle IAM de l'instance. La mise sous DVC a aussi permis d'écarter un dossier
parasite \texttt{\_\_MACOSX} (20 faux fichiers \texttt{.jpg}) présent dans l'archive PBC.

\paragraph{Intégration et livraison continues.} Deux workflows GitHub Actions
accompagnent le dépôt :
\begin{itemize}[leftmargin=*]
  \item \texttt{ci.yml}, à chaque pull request : contrôle du style (\texttt{ruff})
        et tests unitaires (\texttt{pytest}), avec PyTorch en version CPU ;
  \item \texttt{build-bento.yml}, à chaque fusion dans \texttt{main} : construction
        de l'image de l'API, test de fumée (démarrage, refus d'un intrus sur
        \texttt{/login}, refus de \texttt{/predict} sans jeton), puis publication
        dans Amazon ECR, où seules les trois dernières images sont conservées.
\end{itemize}

Depuis le 24 octobre, la chaîne publie l'image de l'API à chaque fusion, et la VM la
récupère directement depuis ECR (\texttt{docker pull}) au lieu de la reconstruire, soit
environ 15 minutes de construction en moins par mise à jour. La mise au point a corrigé
deux problèmes : un conflit de version de \texttt{bentoml} entre \texttt{pyproject.toml} et
\texttt{bentofile.yaml} (la CI installe désormais la version épinglée), et le format d'identité OIDC de l'organisation GitHub, qui inclut des identifiants
numériques de l'organisation et du dépôt (relevé dans CloudTrail) ; la relation
d'approbation du rôle accepte désormais exactement cette identité, sur la seule
branche \texttt{main}.

\paragraph{Supervision.} Prometheus collecte toutes les 15~s les métriques exposées
par l'API ; un tableau de bord Grafana (débit, latence, taux d'erreur, refus
d'authentification, prédictions par classe) est versionné dans le dépôt. Quatre
règles d'alerte sont évaluées en continu :

\begin{table}[H]
\centering
\small
\begin{tabularx}{\linewidth}{@{}l X l@{}}
\toprule
\textbf{Alerte} & \textbf{Condition} & \textbf{Gravité} \\
\midrule
\texttt{APIInjoignable} & Prometheus ne peut plus collecter l'API pendant 1~min & critique \\
\texttt{TauxErreurs5xxEleve} & plus de 5~\% d'erreurs serveur sur 5~min & haute \\
\texttt{LatenceAPIElevee} & p95 de \texttt{/predict} au-delà de 2~s pendant 5~min & moyenne \\
\texttt{RefusAuthentificationNombreux} & plus de 50 refus 401 en 5~min & moyenne \\
\bottomrule
\end{tabularx}
\caption{Règles d'alerte Prometheus de l'API (\texttt{deployment/prometheus/rules/api.yml}).}
\label{tab:alertes-aws}
\end{table}

\paragraph{Exploitation.} La VM s'arrête automatiquement chaque soir à 22~h
(EventBridge Scheduler) et la pile redémarre seule au démarrage suivant. Le disque
a été porté à 80~Go et 4~Go de swap protègent les entraînements d'un manque de
mémoire. Le contexte de construction des images exclut les données et le cache DVC
(\texttt{.dockerignore}) : sans cette exclusion, Docker copiait 12~Go d'images dans
l'image d'entraînement et saturait le disque.

\subsection{Validation de l'API déployée}

Les tests suivants ont été menés sur la VM le 22 octobre 2026, avant le premier
entraînement, puis le 24 octobre après un entraînement de validation du pipeline
(modèle \emph{baseline}, 1\,000 images, 5 époques, F1 macro de test 0,44) promu
\emph{champion} : la chaîne complète, de l'entraînement à la prédiction authentifiée,
fonctionne.

\begin{table}[H]
\centering
\small
\begin{tabularx}{\linewidth}{@{}X c X@{}}
\toprule
\textbf{Requête} & \textbf{Code HTTP} & \textbf{Réponse} \\
\midrule
\texttt{GET /healthz} & 200 & service démarré \\
\texttt{POST /login}, mauvais identifiants & 401 & \texttt{Invalid credentials} \\
\texttt{POST /login}, bons identifiants & 200 & jeton JWT \\
\texttt{POST /predict} sans jeton & 401 & \texttt{Missing authentication token} \\
\texttt{POST /predict}, jeton invalide & 401 & \texttt{Invalid token} \\
\texttt{POST /predict}, jeton valide & 400 & modèle ou alias \emph{champion} absent (attendu) \\
\texttt{GET /models} & 200 & liste vide (attendu) \\
\texttt{GET /metrics} & 200 & compteurs \texttt{blood\_cells\_server\_*} \\
\midrule
\multicolumn{3}{@{}l}{\textit{Après promotion d'un modèle \emph{champion} (24 octobre)}} \\
\texttt{POST /predict}, jeton valide & 200 & classe prédite et probabilités des 10 classes \\
\texttt{GET /models} & 200 & \texttt{baseline-type} \\
\bottomrule
\end{tabularx}
\caption{Tests de l'API sur la VM : l'authentification protège la prédiction, et un modèle promu \emph{champion} est servi.}
\label{tab:tests-api-aws}
\end{table}
, entre sections/mlops et
% sections/conclusion.
% Utilise uniquement des paquets déjà chargés par report.tex :
% tikz (positioning, arrows.meta), booktabs, tabularx, xcolor (couleurs liora*).
% ---------------------------------------------------------------------------

\section{Déploiement sur AWS}
\label{sec:deploiement-aws}

Le service de classification tourne sur une machine virtuelle AWS partagée par
l'équipe, en région Paris (\texttt{eu-west-3}). Toute la chaîne MLOps y est
conteneurisée : suivi des expériences et registre de modèles (MLflow), stockage
des artefacts (MinIO), API de prédiction authentifiée (BentoML), supervision
(Prometheus, Grafana). Les jeux de données sont versionnés avec DVC dans un bucket
S3, et le code est contrôlé puis empaqueté par GitHub Actions. La machine n'expose
aucun port sur Internet : on n'y accède qu'à travers AWS Systems Manager, avec les
droits IAM de chacun.

\subsection{Architecture}

\begin{figure}[H]
\centering
\resizebox{\linewidth}{!}{%
\begin{tikzpicture}[
  font=\sffamily\scriptsize,
  >={Stealth[length=2.2mm]},
  box/.style={draw=lioraink, rounded corners=2pt, fill=white, align=center,
              line width=0.6pt, inner sep=3pt},
  svc/.style={box, minimum width=2.8cm, minimum height=1.5cm, text width=2.6cm},
  ctr/.style={box, minimum width=4.2cm, minimum height=2.0cm, text width=4.0cm},
  ext/.style={box, fill=lioraline!45, minimum height=1.3cm},
  flow/.style={->, draw=lioraink, line width=0.6pt},
  sec/.style={->, draw=liorared, line width=0.8pt},
  lbl/.style={font=\sffamily\tiny, text=lioragrey, align=center, fill=white, inner sep=1pt},
]
% --- Acteurs externes --------------------------------------------------------
\node[ext, minimum width=5.9cm, text width=5.7cm] (team) at (3.35,12.65)
  {\textbf{Postes de l'équipe} (3 membres)\\AWS CLI \texttt{aws login} + MFA,
   plugin Session Manager};
\node[ext, minimum width=5.9cm, text width=5.7cm] (gh) at (9.55,12.65)
  {\textbf{GitHub} \texttt{jul26\_bds\_blood\_cells}\\Actions : CI (ruff, pytest),
   image de l'API testée puis publiée};

% --- Compte AWS -------------------------------------------------------------
\draw[lioraink, line width=0.8pt, rounded corners=4pt] (0,0) rectangle (16,11.2);
\node[anchor=north east, font=\sffamily\small\bfseries, text=lioraink] at (15.85,11.12)
  {AWS \textperiodcentered{} région eu-west-3 (Paris)};

\node[svc] (iam)   at (1.8,9.6)  {\textbf{IAM}\\groupe \texttt{blood-cells-team}\\garde-fous :\\région, suppression};
\node[svc] (ssm)   at (4.9,9.6)  {\textbf{Systems Manager}\\Session Manager\\(shell, tunnels)};
\node[svc] (ecr)   at (8.0,9.6)  {\textbf{ECR}\\registre d'images\\(image de l'API)};
\node[svc] (s3)    at (11.1,9.6) {\textbf{S3}\\jeux de données\\versionnés (DVC)};
\node[svc] (sched) at (14.2,9.6) {\textbf{EventBridge Scheduler}\\arrêt de la VM à 22 h};

% --- Machine virtuelle --------------------------------------------------------
\draw[liorared, line width=0.9pt, rounded corners=4pt] (0.3,0.3) rectangle (15.7,8.0);
\node[anchor=north west, align=left, text=lioraink] at (0.45,7.95)
  {\textbf{EC2 \texttt{blood-cells-vm}} \textperiodcentered{} m7i-flex.large (2 vCPU, 8 Go, 80 Go)
   \textperiodcentered{} Ubuntu 24.04 \textperiodcentered{} \textcolor{liorared}{aucun port entrant}
   \textperiodcentered{} rôle \texttt{blood-cells-ec2-role}};

\draw[lioragrey, dashed, rounded corners=3pt] (0.6,0.6) rectangle (15.4,7.1);
\node[anchor=north west, text=lioragrey] at (0.7,7.05) {Docker Compose};

\node[ctr] (api)   at (3.1,5.2)  {\textbf{API BentoML} (8000)\\\texttt{/login} (JWT),
  \texttt{/predict}, \texttt{/models}, \texttt{/metrics}};
\node[ctr] (mlf)   at (7.9,5.2)  {\textbf{MLflow} (5001)\\suivi des expériences,\\registre, alias \emph{champion}};
\node[ctr] (minio) at (12.7,5.2) {\textbf{MinIO} (9000, 9001)\\stockage S3 des artefacts MLflow};
\node[ctr] (prom)  at (3.1,2.45) {\textbf{Prometheus} (9090)\\collecte des métriques,\\4 règles d'alerte};
\node[ctr] (train) at (7.9,2.45) {\textbf{Entraînement}\\\texttt{src.experiments.run}\\lit \texttt{data/} (PBC, LMU)};
\node[ctr] (graf)  at (12.7,2.45) {\textbf{Grafana} (9091)\\tableau de bord de l'API};

% --- Flux -------------------------------------------------------------------
\draw[sec]  (team.south -| iam)  -- node[lbl, pos=0.45] {authentification} (iam.north);
\draw[sec]  (team.south -| ssm)  -- node[lbl, pos=0.45] {session chiffrée} (ssm.north);
\draw[sec]  (gh.south -| ecr)    -- node[lbl, pos=0.45] {OIDC, \texttt{docker push}} (ecr.north);

\draw[sec]  (ssm.south)   -- node[lbl] {shell et\\tunnels de ports} (ssm.south |- 0,8.0);
\draw[flow] (ecr.south)   -- node[lbl] {\texttt{docker pull}} (ecr.south |- 0,8.0);
\draw[flow] (s3.south)    -- node[lbl] {\texttt{dvc pull}} (s3.south |- 0,8.0);
\draw[flow] (sched.south) -- node[lbl] {arrêt} (sched.south |- 0,8.0);

\draw[flow] (api)   -- node[lbl, above=1pt] {charge le\\\emph{champion}} (mlf);
\draw[flow] (mlf)   -- node[lbl, above=1pt] {artefacts} (minio);
\draw[flow] (train) -- node[lbl] {runs, modèles} (mlf);
\draw[flow] (prom)  -- node[lbl] {lit \texttt{/metrics}} (api);
\draw[flow] (graf.south) -- ++(0,-0.45) -| node[lbl, pos=0.25] {requêtes} (prom.south);
\end{tikzpicture}}
\caption{Architecture déployée. En rouge, les accès contrôlés par IAM : personne
n'ouvre de port réseau, tout passe par AWS Systems Manager ou par la fédération
OIDC de GitHub. En pointillés, les services de la pile Docker Compose.}
\label{fig:architecture-aws}
\end{figure}

\subsection{Services et accès}

\begin{table}[H]
\centering
\small
\begin{tabularx}{\linewidth}{@{}l l X@{}}
\toprule
\textbf{Service} & \textbf{Port (VM)} & \textbf{Rôle} \\
\midrule
API BentoML & 8000 & prédiction ; \texttt{/login} délivre un jeton JWT valable 1 h, exigé par \texttt{/predict} \\
MLflow      & 5001 & suivi des entraînements et registre ; l'API sert la version d'alias \emph{champion} \\
MinIO       & 9000, 9001 & stockage objet des artefacts MLflow (API S3, console) \\
Prometheus  & 9090 & collecte toutes les 15 s les métriques exposées par l'API \\
Grafana     & 9091 & tableaux de bord de supervision \\
\bottomrule
\end{tabularx}
\caption{Services de la pile Docker Compose. Les ports ne sont joignables que
depuis la VM ou par un tunnel Session Manager ouvert depuis le poste d'un membre.}
\label{tab:services-aws}
\end{table}

\subsection{Sécurité}

\begin{itemize}[leftmargin=*]
  \item \textbf{Identités.} Chaque membre a son utilisateur IAM, membre du groupe
        \texttt{blood-cells-team}, avec MFA. La CLI se connecte par
        \texttt{aws login} : aucune clé d'accès permanente n'est stockée sur les postes.
  \item \textbf{Garde-fous.} Une politique de groupe refuse toute action hors
        des régions Paris et \texttt{us-east-1} (services globaux), ainsi que la
        suppression de la VM, de son disque et du bucket (étiquette \texttt{Protected}).
  \item \textbf{Réseau.} Le groupe de sécurité de la VM n'a aucune règle entrante ;
        ni SSH ni clé SSH. Shell et tunnels passent par Session Manager.
  \item \textbf{Secrets.} Mots de passe MinIO et API, clé de signature JWT :
        générés aléatoirement dans un fichier \texttt{.env} présent uniquement sur
        la VM (droits 600, ignoré par Git). La VM lit le dépôt avec un jeton
        GitHub en lecture seule. GitHub Actions publie dans ECR par fédération
        OIDC : aucun secret AWS n'est stocké dans GitHub, et seule la branche
        \texttt{main} peut endosser le rôle de publication.
  \item \textbf{Coût.} La VM est arrêtée automatiquement chaque soir à 22 h, et des
        alertes de facturation sont actives.
\end{itemize}

\subsection{Données, intégration continue et supervision}

\paragraph{Données versionnées.} Les deux jeux d'images (PBC et LMU) sont suivis par
DVC : le dépôt Git ne contient que deux petits fichiers \texttt{.dvc}, qui décrivent
le contenu exact de chaque dossier, tandis que les images sont stockées dans
\texttt{s3://blood-cells-123456789012/dvc} (35\,451 objets, 12,0~Go, les fichiers
identiques n'étant stockés qu'une fois). Chaque commit référence donc la version
précise des données ; sur la VM, \texttt{dvc pull} les restaure en quelques secondes
grâce au rôle IAM de l'instance. La mise sous DVC a aussi permis d'écarter un dossier
parasite \texttt{\_\_MACOSX} (20 faux fichiers \texttt{.jpg}) présent dans l'archive PBC.

\paragraph{Intégration et livraison continues.} Deux workflows GitHub Actions
accompagnent le dépôt :
\begin{itemize}[leftmargin=*]
  \item \texttt{ci.yml}, à chaque pull request : contrôle du style (\texttt{ruff})
        et tests unitaires (\texttt{pytest}), avec PyTorch en version CPU ;
  \item \texttt{build-bento.yml}, à chaque fusion dans \texttt{main} : construction
        de l'image de l'API, test de fumée (démarrage, refus d'un intrus sur
        \texttt{/login}, refus de \texttt{/predict} sans jeton), puis publication
        dans Amazon ECR, où seules les trois dernières images sont conservées.
\end{itemize}

Depuis le 24 octobre, la chaîne publie l'image de l'API à chaque fusion, et la VM la
récupère directement depuis ECR (\texttt{docker pull}) au lieu de la reconstruire, soit
environ 15 minutes de construction en moins par mise à jour. La mise au point a corrigé
deux problèmes : un conflit de version de \texttt{bentoml} entre \texttt{pyproject.toml} et
\texttt{bentofile.yaml} (la CI installe désormais la version épinglée), et le format d'identité OIDC de l'organisation GitHub, qui inclut des identifiants
numériques de l'organisation et du dépôt (relevé dans CloudTrail) ; la relation
d'approbation du rôle accepte désormais exactement cette identité, sur la seule
branche \texttt{main}.

\paragraph{Supervision.} Prometheus collecte toutes les 15~s les métriques exposées
par l'API ; un tableau de bord Grafana (débit, latence, taux d'erreur, refus
d'authentification, prédictions par classe) est versionné dans le dépôt. Quatre
règles d'alerte sont évaluées en continu :

\begin{table}[H]
\centering
\small
\begin{tabularx}{\linewidth}{@{}l X l@{}}
\toprule
\textbf{Alerte} & \textbf{Condition} & \textbf{Gravité} \\
\midrule
\texttt{APIInjoignable} & Prometheus ne peut plus collecter l'API pendant 1~min & critique \\
\texttt{TauxErreurs5xxEleve} & plus de 5~\% d'erreurs serveur sur 5~min & haute \\
\texttt{LatenceAPIElevee} & p95 de \texttt{/predict} au-delà de 2~s pendant 5~min & moyenne \\
\texttt{RefusAuthentificationNombreux} & plus de 50 refus 401 en 5~min & moyenne \\
\bottomrule
\end{tabularx}
\caption{Règles d'alerte Prometheus de l'API (\texttt{deployment/prometheus/rules/api.yml}).}
\label{tab:alertes-aws}
\end{table}

\paragraph{Exploitation.} La VM s'arrête automatiquement chaque soir à 22~h
(EventBridge Scheduler) et la pile redémarre seule au démarrage suivant. Le disque
a été porté à 80~Go et 4~Go de swap protègent les entraînements d'un manque de
mémoire. Le contexte de construction des images exclut les données et le cache DVC
(\texttt{.dockerignore}) : sans cette exclusion, Docker copiait 12~Go d'images dans
l'image d'entraînement et saturait le disque.

\subsection{Validation de l'API déployée}

Les tests suivants ont été menés sur la VM le 22 octobre 2026, avant le premier
entraînement, puis le 24 octobre après un entraînement de validation du pipeline
(modèle \emph{baseline}, 1\,000 images, 5 époques, F1 macro de test 0,44) promu
\emph{champion} : la chaîne complète, de l'entraînement à la prédiction authentifiée,
fonctionne.

\begin{table}[H]
\centering
\small
\begin{tabularx}{\linewidth}{@{}X c X@{}}
\toprule
\textbf{Requête} & \textbf{Code HTTP} & \textbf{Réponse} \\
\midrule
\texttt{GET /healthz} & 200 & service démarré \\
\texttt{POST /login}, mauvais identifiants & 401 & \texttt{Invalid credentials} \\
\texttt{POST /login}, bons identifiants & 200 & jeton JWT \\
\texttt{POST /predict} sans jeton & 401 & \texttt{Missing authentication token} \\
\texttt{POST /predict}, jeton invalide & 401 & \texttt{Invalid token} \\
\texttt{POST /predict}, jeton valide & 400 & modèle ou alias \emph{champion} absent (attendu) \\
\texttt{GET /models} & 200 & liste vide (attendu) \\
\texttt{GET /metrics} & 200 & compteurs \texttt{blood\_cells\_server\_*} \\
\midrule
\multicolumn{3}{@{}l}{\textit{Après promotion d'un modèle \emph{champion} (24 octobre)}} \\
\texttt{POST /predict}, jeton valide & 200 & classe prédite et probabilités des 10 classes \\
\texttt{GET /models} & 200 & \texttt{baseline-type} \\
\bottomrule
\end{tabularx}
\caption{Tests de l'API sur la VM : l'authentification protège la prédiction, et un modèle promu \emph{champion} est servi.}
\label{tab:tests-api-aws}
\end{table}
, entre sections/mlops et
% sections/conclusion.
% Utilise uniquement des paquets déjà chargés par report.tex :
% tikz (positioning, arrows.meta), booktabs, tabularx, xcolor (couleurs liora*).
% ---------------------------------------------------------------------------

\section{Déploiement sur AWS}
\label{sec:deploiement-aws}

Le service de classification tourne sur une machine virtuelle AWS partagée par
l'équipe, en région Paris (\texttt{eu-west-3}). Toute la chaîne MLOps y est
conteneurisée : suivi des expériences et registre de modèles (MLflow), stockage
des artefacts (MinIO), API de prédiction authentifiée (BentoML), supervision
(Prometheus, Grafana). Les jeux de données sont versionnés avec DVC dans un bucket
S3, et le code est contrôlé puis empaqueté par GitHub Actions. La machine n'expose
aucun port sur Internet : on n'y accède qu'à travers AWS Systems Manager, avec les
droits IAM de chacun.

\subsection{Architecture}

\begin{figure}[H]
\centering
\resizebox{\linewidth}{!}{%
\begin{tikzpicture}[
  font=\sffamily\scriptsize,
  >={Stealth[length=2.2mm]},
  box/.style={draw=lioraink, rounded corners=2pt, fill=white, align=center,
              line width=0.6pt, inner sep=3pt},
  svc/.style={box, minimum width=2.8cm, minimum height=1.5cm, text width=2.6cm},
  ctr/.style={box, minimum width=4.2cm, minimum height=2.0cm, text width=4.0cm},
  ext/.style={box, fill=lioraline!45, minimum height=1.3cm},
  flow/.style={->, draw=lioraink, line width=0.6pt},
  sec/.style={->, draw=liorared, line width=0.8pt},
  lbl/.style={font=\sffamily\tiny, text=lioragrey, align=center, fill=white, inner sep=1pt},
]
% --- Acteurs externes --------------------------------------------------------
\node[ext, minimum width=5.9cm, text width=5.7cm] (team) at (3.35,12.65)
  {\textbf{Postes de l'équipe} (3 membres)\\AWS CLI \texttt{aws login} + MFA,
   plugin Session Manager};
\node[ext, minimum width=5.9cm, text width=5.7cm] (gh) at (9.55,12.65)
  {\textbf{GitHub} \texttt{jul26\_bds\_blood\_cells}\\Actions : CI (ruff, pytest),
   image de l'API testée puis publiée};

% --- Compte AWS -------------------------------------------------------------
\draw[lioraink, line width=0.8pt, rounded corners=4pt] (0,0) rectangle (16,11.2);
\node[anchor=north east, font=\sffamily\small\bfseries, text=lioraink] at (15.85,11.12)
  {AWS \textperiodcentered{} région eu-west-3 (Paris)};

\node[svc] (iam)   at (1.8,9.6)  {\textbf{IAM}\\groupe \texttt{blood-cells-team}\\garde-fous :\\région, suppression};
\node[svc] (ssm)   at (4.9,9.6)  {\textbf{Systems Manager}\\Session Manager\\(shell, tunnels)};
\node[svc] (ecr)   at (8.0,9.6)  {\textbf{ECR}\\registre d'images\\(image de l'API)};
\node[svc] (s3)    at (11.1,9.6) {\textbf{S3}\\jeux de données\\versionnés (DVC)};
\node[svc] (sched) at (14.2,9.6) {\textbf{EventBridge Scheduler}\\arrêt de la VM à 22 h};

% --- Machine virtuelle --------------------------------------------------------
\draw[liorared, line width=0.9pt, rounded corners=4pt] (0.3,0.3) rectangle (15.7,8.0);
\node[anchor=north west, align=left, text=lioraink] at (0.45,7.95)
  {\textbf{EC2 \texttt{blood-cells-vm}} \textperiodcentered{} m7i-flex.large (2 vCPU, 8 Go, 80 Go)
   \textperiodcentered{} Ubuntu 24.04 \textperiodcentered{} \textcolor{liorared}{aucun port entrant}
   \textperiodcentered{} rôle \texttt{blood-cells-ec2-role}};

\draw[lioragrey, dashed, rounded corners=3pt] (0.6,0.6) rectangle (15.4,7.1);
\node[anchor=north west, text=lioragrey] at (0.7,7.05) {Docker Compose};

\node[ctr] (api)   at (3.1,5.2)  {\textbf{API BentoML} (8000)\\\texttt{/login} (JWT),
  \texttt{/predict}, \texttt{/models}, \texttt{/metrics}};
\node[ctr] (mlf)   at (7.9,5.2)  {\textbf{MLflow} (5001)\\suivi des expériences,\\registre, alias \emph{champion}};
\node[ctr] (minio) at (12.7,5.2) {\textbf{MinIO} (9000, 9001)\\stockage S3 des artefacts MLflow};
\node[ctr] (prom)  at (3.1,2.45) {\textbf{Prometheus} (9090)\\collecte des métriques,\\4 règles d'alerte};
\node[ctr] (train) at (7.9,2.45) {\textbf{Entraînement}\\\texttt{src.experiments.run}\\lit \texttt{data/} (PBC, LMU)};
\node[ctr] (graf)  at (12.7,2.45) {\textbf{Grafana} (9091)\\tableau de bord de l'API};

% --- Flux -------------------------------------------------------------------
\draw[sec]  (team.south -| iam)  -- node[lbl, pos=0.45] {authentification} (iam.north);
\draw[sec]  (team.south -| ssm)  -- node[lbl, pos=0.45] {session chiffrée} (ssm.north);
\draw[sec]  (gh.south -| ecr)    -- node[lbl, pos=0.45] {OIDC, \texttt{docker push}} (ecr.north);

\draw[sec]  (ssm.south)   -- node[lbl] {shell et\\tunnels de ports} (ssm.south |- 0,8.0);
\draw[flow] (ecr.south)   -- node[lbl] {\texttt{docker pull}} (ecr.south |- 0,8.0);
\draw[flow] (s3.south)    -- node[lbl] {\texttt{dvc pull}} (s3.south |- 0,8.0);
\draw[flow] (sched.south) -- node[lbl] {arrêt} (sched.south |- 0,8.0);

\draw[flow] (api)   -- node[lbl, above=1pt] {charge le\\\emph{champion}} (mlf);
\draw[flow] (mlf)   -- node[lbl, above=1pt] {artefacts} (minio);
\draw[flow] (train) -- node[lbl] {runs, modèles} (mlf);
\draw[flow] (prom)  -- node[lbl] {lit \texttt{/metrics}} (api);
\draw[flow] (graf.south) -- ++(0,-0.45) -| node[lbl, pos=0.25] {requêtes} (prom.south);
\end{tikzpicture}}
\caption{Architecture déployée. En rouge, les accès contrôlés par IAM : personne
n'ouvre de port réseau, tout passe par AWS Systems Manager ou par la fédération
OIDC de GitHub. En pointillés, les services de la pile Docker Compose.}
\label{fig:architecture-aws}
\end{figure}

\subsection{Services et accès}

\begin{table}[H]
\centering
\small
\begin{tabularx}{\linewidth}{@{}l l X@{}}
\toprule
\textbf{Service} & \textbf{Port (VM)} & \textbf{Rôle} \\
\midrule
API BentoML & 8000 & prédiction ; \texttt{/login} délivre un jeton JWT valable 1 h, exigé par \texttt{/predict} \\
MLflow      & 5001 & suivi des entraînements et registre ; l'API sert la version d'alias \emph{champion} \\
MinIO       & 9000, 9001 & stockage objet des artefacts MLflow (API S3, console) \\
Prometheus  & 9090 & collecte toutes les 15 s les métriques exposées par l'API \\
Grafana     & 9091 & tableaux de bord de supervision \\
\bottomrule
\end{tabularx}
\caption{Services de la pile Docker Compose. Les ports ne sont joignables que
depuis la VM ou par un tunnel Session Manager ouvert depuis le poste d'un membre.}
\label{tab:services-aws}
\end{table}

\subsection{Sécurité}

\begin{itemize}[leftmargin=*]
  \item \textbf{Identités.} Chaque membre a son utilisateur IAM, membre du groupe
        \texttt{blood-cells-team}, avec MFA. La CLI se connecte par
        \texttt{aws login} : aucune clé d'accès permanente n'est stockée sur les postes.
  \item \textbf{Garde-fous.} Une politique de groupe refuse toute action hors
        des régions Paris et \texttt{us-east-1} (services globaux), ainsi que la
        suppression de la VM, de son disque et du bucket (étiquette \texttt{Protected}).
  \item \textbf{Réseau.} Le groupe de sécurité de la VM n'a aucune règle entrante ;
        ni SSH ni clé SSH. Shell et tunnels passent par Session Manager.
  \item \textbf{Secrets.} Mots de passe MinIO et API, clé de signature JWT :
        générés aléatoirement dans un fichier \texttt{.env} présent uniquement sur
        la VM (droits 600, ignoré par Git). La VM lit le dépôt avec un jeton
        GitHub en lecture seule. GitHub Actions publie dans ECR par fédération
        OIDC : aucun secret AWS n'est stocké dans GitHub, et seule la branche
        \texttt{main} peut endosser le rôle de publication.
  \item \textbf{Coût.} La VM est arrêtée automatiquement chaque soir à 22 h, et des
        alertes de facturation sont actives.
\end{itemize}

\subsection{Données, intégration continue et supervision}

\paragraph{Données versionnées.} Les deux jeux d'images (PBC et LMU) sont suivis par
DVC : le dépôt Git ne contient que deux petits fichiers \texttt{.dvc}, qui décrivent
le contenu exact de chaque dossier, tandis que les images sont stockées dans
\texttt{s3://blood-cells-123456789012/dvc} (35\,451 objets, 12,0~Go, les fichiers
identiques n'étant stockés qu'une fois). Chaque commit référence donc la version
précise des données ; sur la VM, \texttt{dvc pull} les restaure en quelques secondes
grâce au rôle IAM de l'instance. La mise sous DVC a aussi permis d'écarter un dossier
parasite \texttt{\_\_MACOSX} (20 faux fichiers \texttt{.jpg}) présent dans l'archive PBC.

\paragraph{Intégration et livraison continues.} Deux workflows GitHub Actions
accompagnent le dépôt :
\begin{itemize}[leftmargin=*]
  \item \texttt{ci.yml}, à chaque pull request : contrôle du style (\texttt{ruff})
        et tests unitaires (\texttt{pytest}), avec PyTorch en version CPU ;
  \item \texttt{build-bento.yml}, à chaque fusion dans \texttt{main} : construction
        de l'image de l'API, test de fumée (démarrage, refus d'un intrus sur
        \texttt{/login}, refus de \texttt{/predict} sans jeton), puis publication
        dans Amazon ECR, où seules les trois dernières images sont conservées.
\end{itemize}

Depuis le 24 octobre, la chaîne publie l'image de l'API à chaque fusion, et la VM la
récupère directement depuis ECR (\texttt{docker pull}) au lieu de la reconstruire, soit
environ 15 minutes de construction en moins par mise à jour. La mise au point a corrigé
deux problèmes : un conflit de version de \texttt{bentoml} entre \texttt{pyproject.toml} et
\texttt{bentofile.yaml} (la CI installe désormais la version épinglée), et le format d'identité OIDC de l'organisation GitHub, qui inclut des identifiants
numériques de l'organisation et du dépôt (relevé dans CloudTrail) ; la relation
d'approbation du rôle accepte désormais exactement cette identité, sur la seule
branche \texttt{main}.

\paragraph{Supervision.} Prometheus collecte toutes les 15~s les métriques exposées
par l'API ; un tableau de bord Grafana (débit, latence, taux d'erreur, refus
d'authentification, prédictions par classe) est versionné dans le dépôt. Quatre
règles d'alerte sont évaluées en continu :

\begin{table}[H]
\centering
\small
\begin{tabularx}{\linewidth}{@{}l X l@{}}
\toprule
\textbf{Alerte} & \textbf{Condition} & \textbf{Gravité} \\
\midrule
\texttt{APIInjoignable} & Prometheus ne peut plus collecter l'API pendant 1~min & critique \\
\texttt{TauxErreurs5xxEleve} & plus de 5~\% d'erreurs serveur sur 5~min & haute \\
\texttt{LatenceAPIElevee} & p95 de \texttt{/predict} au-delà de 2~s pendant 5~min & moyenne \\
\texttt{RefusAuthentificationNombreux} & plus de 50 refus 401 en 5~min & moyenne \\
\bottomrule
\end{tabularx}
\caption{Règles d'alerte Prometheus de l'API (\texttt{deployment/prometheus/rules/api.yml}).}
\label{tab:alertes-aws}
\end{table}

\paragraph{Exploitation.} La VM s'arrête automatiquement chaque soir à 22~h
(EventBridge Scheduler) et la pile redémarre seule au démarrage suivant. Le disque
a été porté à 80~Go et 4~Go de swap protègent les entraînements d'un manque de
mémoire. Le contexte de construction des images exclut les données et le cache DVC
(\texttt{.dockerignore}) : sans cette exclusion, Docker copiait 12~Go d'images dans
l'image d'entraînement et saturait le disque.

\subsection{Validation de l'API déployée}

Les tests suivants ont été menés sur la VM le 22 octobre 2026, avant le premier
entraînement, puis le 24 octobre après un entraînement de validation du pipeline
(modèle \emph{baseline}, 1\,000 images, 5 époques, F1 macro de test 0,44) promu
\emph{champion} : la chaîne complète, de l'entraînement à la prédiction authentifiée,
fonctionne.

\begin{table}[H]
\centering
\small
\begin{tabularx}{\linewidth}{@{}X c X@{}}
\toprule
\textbf{Requête} & \textbf{Code HTTP} & \textbf{Réponse} \\
\midrule
\texttt{GET /healthz} & 200 & service démarré \\
\texttt{POST /login}, mauvais identifiants & 401 & \texttt{Invalid credentials} \\
\texttt{POST /login}, bons identifiants & 200 & jeton JWT \\
\texttt{POST /predict} sans jeton & 401 & \texttt{Missing authentication token} \\
\texttt{POST /predict}, jeton invalide & 401 & \texttt{Invalid token} \\
\texttt{POST /predict}, jeton valide & 400 & modèle ou alias \emph{champion} absent (attendu) \\
\texttt{GET /models} & 200 & liste vide (attendu) \\
\texttt{GET /metrics} & 200 & compteurs \texttt{blood\_cells\_server\_*} \\
\midrule
\multicolumn{3}{@{}l}{\textit{Après promotion d'un modèle \emph{champion} (24 octobre)}} \\
\texttt{POST /predict}, jeton valide & 200 & classe prédite et probabilités des 10 classes \\
\texttt{GET /models} & 200 & \texttt{baseline-type} \\
\bottomrule
\end{tabularx}
\caption{Tests de l'API sur la VM : l'authentification protège la prédiction, et un modèle promu \emph{champion} est servi.}
\label{tab:tests-api-aws}
\end{table}
, entre sections/mlops et
% sections/conclusion.
% Utilise uniquement des paquets déjà chargés par report.tex :
% tikz (positioning, arrows.meta), booktabs, tabularx, xcolor (couleurs liora*).
% ---------------------------------------------------------------------------

\section{Déploiement sur AWS}
\label{sec:deploiement-aws}

Le service de classification tourne sur une machine virtuelle AWS partagée par
l'équipe, en région Paris (\texttt{eu-west-3}). Toute la chaîne MLOps y est
conteneurisée : suivi des expériences et registre de modèles (MLflow), stockage
des artefacts (MinIO), API de prédiction authentifiée (BentoML), supervision
(Prometheus, Grafana). Les jeux de données sont versionnés avec DVC dans un bucket
S3, et le code est contrôlé puis empaqueté par GitHub Actions. La machine n'expose
aucun port sur Internet : on n'y accède qu'à travers AWS Systems Manager, avec les
droits IAM de chacun.

\subsection{Architecture}

\begin{figure}[H]
\centering
\resizebox{\linewidth}{!}{%
\begin{tikzpicture}[
  font=\sffamily\scriptsize,
  >={Stealth[length=2.2mm]},
  box/.style={draw=lioraink, rounded corners=2pt, fill=white, align=center,
              line width=0.6pt, inner sep=3pt},
  svc/.style={box, minimum width=2.8cm, minimum height=1.5cm, text width=2.6cm},
  ctr/.style={box, minimum width=4.2cm, minimum height=2.0cm, text width=4.0cm},
  ext/.style={box, fill=lioraline!45, minimum height=1.3cm},
  flow/.style={->, draw=lioraink, line width=0.6pt},
  sec/.style={->, draw=liorared, line width=0.8pt},
  lbl/.style={font=\sffamily\tiny, text=lioragrey, align=center, fill=white, inner sep=1pt},
]
% --- Acteurs externes --------------------------------------------------------
\node[ext, minimum width=5.9cm, text width=5.7cm] (team) at (3.35,12.65)
  {\textbf{Postes de l'équipe} (3 membres)\\AWS CLI \texttt{aws login} + MFA,
   plugin Session Manager};
\node[ext, minimum width=5.9cm, text width=5.7cm] (gh) at (9.55,12.65)
  {\textbf{GitHub} \texttt{jul26\_bds\_blood\_cells}\\Actions : CI (ruff, pytest),
   image de l'API testée puis publiée};

% --- Compte AWS -------------------------------------------------------------
\draw[lioraink, line width=0.8pt, rounded corners=4pt] (0,0) rectangle (16,11.2);
\node[anchor=north east, font=\sffamily\small\bfseries, text=lioraink] at (15.85,11.12)
  {AWS \textperiodcentered{} région eu-west-3 (Paris)};

\node[svc] (iam)   at (1.8,9.6)  {\textbf{IAM}\\groupe \texttt{blood-cells-team}\\garde-fous :\\région, suppression};
\node[svc] (ssm)   at (4.9,9.6)  {\textbf{Systems Manager}\\Session Manager\\(shell, tunnels)};
\node[svc] (ecr)   at (8.0,9.6)  {\textbf{ECR}\\registre d'images\\(image de l'API)};
\node[svc] (s3)    at (11.1,9.6) {\textbf{S3}\\jeux de données\\versionnés (DVC)};
\node[svc] (sched) at (14.2,9.6) {\textbf{EventBridge Scheduler}\\arrêt de la VM à 22 h};

% --- Machine virtuelle --------------------------------------------------------
\draw[liorared, line width=0.9pt, rounded corners=4pt] (0.3,0.3) rectangle (15.7,8.0);
\node[anchor=north west, align=left, text=lioraink] at (0.45,7.95)
  {\textbf{EC2 \texttt{blood-cells-vm}} \textperiodcentered{} m7i-flex.large (2 vCPU, 8 Go, 80 Go)
   \textperiodcentered{} Ubuntu 24.04 \textperiodcentered{} \textcolor{liorared}{aucun port entrant}
   \textperiodcentered{} rôle \texttt{blood-cells-ec2-role}};

\draw[lioragrey, dashed, rounded corners=3pt] (0.6,0.6) rectangle (15.4,7.1);
\node[anchor=north west, text=lioragrey] at (0.7,7.05) {Docker Compose};

\node[ctr] (api)   at (3.1,5.2)  {\textbf{API BentoML} (8000)\\\texttt{/login} (JWT),
  \texttt{/predict}, \texttt{/models}, \texttt{/metrics}};
\node[ctr] (mlf)   at (7.9,5.2)  {\textbf{MLflow} (5001)\\suivi des expériences,\\registre, alias \emph{champion}};
\node[ctr] (minio) at (12.7,5.2) {\textbf{MinIO} (9000, 9001)\\stockage S3 des artefacts MLflow};
\node[ctr] (prom)  at (3.1,2.45) {\textbf{Prometheus} (9090)\\collecte des métriques,\\4 règles d'alerte};
\node[ctr] (train) at (7.9,2.45) {\textbf{Entraînement}\\\texttt{src.experiments.run}\\lit \texttt{data/} (PBC, LMU)};
\node[ctr] (graf)  at (12.7,2.45) {\textbf{Grafana} (9091)\\tableau de bord de l'API};

% --- Flux -------------------------------------------------------------------
\draw[sec]  (team.south -| iam)  -- node[lbl, pos=0.45] {authentification} (iam.north);
\draw[sec]  (team.south -| ssm)  -- node[lbl, pos=0.45] {session chiffrée} (ssm.north);
\draw[sec]  (gh.south -| ecr)    -- node[lbl, pos=0.45] {OIDC, \texttt{docker push}} (ecr.north);

\draw[sec]  (ssm.south)   -- node[lbl] {shell et\\tunnels de ports} (ssm.south |- 0,8.0);
\draw[flow] (ecr.south)   -- node[lbl] {\texttt{docker pull}} (ecr.south |- 0,8.0);
\draw[flow] (s3.south)    -- node[lbl] {\texttt{dvc pull}} (s3.south |- 0,8.0);
\draw[flow] (sched.south) -- node[lbl] {arrêt} (sched.south |- 0,8.0);

\draw[flow] (api)   -- node[lbl, above=1pt] {charge le\\\emph{champion}} (mlf);
\draw[flow] (mlf)   -- node[lbl, above=1pt] {artefacts} (minio);
\draw[flow] (train) -- node[lbl] {runs, modèles} (mlf);
\draw[flow] (prom)  -- node[lbl] {lit \texttt{/metrics}} (api);
\draw[flow] (graf.south) -- ++(0,-0.45) -| node[lbl, pos=0.25] {requêtes} (prom.south);
\end{tikzpicture}}
\caption{Architecture déployée. En rouge, les accès contrôlés par IAM : personne
n'ouvre de port réseau, tout passe par AWS Systems Manager ou par la fédération
OIDC de GitHub. En pointillés, les services de la pile Docker Compose.}
\label{fig:architecture-aws}
\end{figure}

\subsection{Services et accès}

\begin{table}[H]
\centering
\small
\begin{tabularx}{\linewidth}{@{}l l X@{}}
\toprule
\textbf{Service} & \textbf{Port (VM)} & \textbf{Rôle} \\
\midrule
API BentoML & 8000 & prédiction ; \texttt{/login} délivre un jeton JWT valable 1 h, exigé par \texttt{/predict} \\
MLflow      & 5001 & suivi des entraînements et registre ; l'API sert la version d'alias \emph{champion} \\
MinIO       & 9000, 9001 & stockage objet des artefacts MLflow (API S3, console) \\
Prometheus  & 9090 & collecte toutes les 15 s les métriques exposées par l'API \\
Grafana     & 9091 & tableaux de bord de supervision \\
\bottomrule
\end{tabularx}
\caption{Services de la pile Docker Compose. Les ports ne sont joignables que
depuis la VM ou par un tunnel Session Manager ouvert depuis le poste d'un membre.}
\label{tab:services-aws}
\end{table}

\subsection{Sécurité}

\begin{itemize}[leftmargin=*]
  \item \textbf{Identités.} Chaque membre a son utilisateur IAM, membre du groupe
        \texttt{blood-cells-team}, avec MFA. La CLI se connecte par
        \texttt{aws login} : aucune clé d'accès permanente n'est stockée sur les postes.
  \item \textbf{Garde-fous.} Une politique de groupe refuse toute action hors
        des régions Paris et \texttt{us-east-1} (services globaux), ainsi que la
        suppression de la VM, de son disque et du bucket (étiquette \texttt{Protected}).
  \item \textbf{Réseau.} Le groupe de sécurité de la VM n'a aucune règle entrante ;
        ni SSH ni clé SSH. Shell et tunnels passent par Session Manager.
  \item \textbf{Secrets.} Mots de passe MinIO et API, clé de signature JWT :
        générés aléatoirement dans un fichier \texttt{.env} présent uniquement sur
        la VM (droits 600, ignoré par Git). La VM lit le dépôt avec un jeton
        GitHub en lecture seule. GitHub Actions publie dans ECR par fédération
        OIDC : aucun secret AWS n'est stocké dans GitHub, et seule la branche
        \texttt{main} peut endosser le rôle de publication.
  \item \textbf{Coût.} La VM est arrêtée automatiquement chaque soir à 22 h, et des
        alertes de facturation sont actives.
\end{itemize}

\subsection{Données, intégration continue et supervision}

\paragraph{Données versionnées.} Les deux jeux d'images (PBC et LMU) sont suivis par
DVC : le dépôt Git ne contient que deux petits fichiers \texttt{.dvc}, qui décrivent
le contenu exact de chaque dossier, tandis que les images sont stockées dans
\texttt{s3://blood-cells-123456789012/dvc} (35\,451 objets, 12,0~Go, les fichiers
identiques n'étant stockés qu'une fois). Chaque commit référence donc la version
précise des données ; sur la VM, \texttt{dvc pull} les restaure en quelques secondes
grâce au rôle IAM de l'instance. La mise sous DVC a aussi permis d'écarter un dossier
parasite \texttt{\_\_MACOSX} (20 faux fichiers \texttt{.jpg}) présent dans l'archive PBC.

\paragraph{Intégration et livraison continues.} Deux workflows GitHub Actions
accompagnent le dépôt :
\begin{itemize}[leftmargin=*]
  \item \texttt{ci.yml}, à chaque pull request : contrôle du style (\texttt{ruff})
        et tests unitaires (\texttt{pytest}), avec PyTorch en version CPU ;
  \item \texttt{build-bento.yml}, à chaque fusion dans \texttt{main} : construction
        de l'image de l'API, test de fumée (démarrage, refus d'un intrus sur
        \texttt{/login}, refus de \texttt{/predict} sans jeton), puis publication
        dans Amazon ECR, où seules les trois dernières images sont conservées.
\end{itemize}

Depuis le 24 octobre, la chaîne publie l'image de l'API à chaque fusion, et la VM la
récupère directement depuis ECR (\texttt{docker pull}) au lieu de la reconstruire, soit
environ 15 minutes de construction en moins par mise à jour. La mise au point a corrigé
deux problèmes : un conflit de version de \texttt{bentoml} entre \texttt{pyproject.toml} et
\texttt{bentofile.yaml} (la CI installe désormais la version épinglée), et le format d'identité OIDC de l'organisation GitHub, qui inclut des identifiants
numériques de l'organisation et du dépôt (relevé dans CloudTrail) ; la relation
d'approbation du rôle accepte désormais exactement cette identité, sur la seule
branche \texttt{main}.

\paragraph{Supervision.} Prometheus collecte toutes les 15~s les métriques exposées
par l'API ; un tableau de bord Grafana (débit, latence, taux d'erreur, refus
d'authentification, prédictions par classe) est versionné dans le dépôt. Quatre
règles d'alerte sont évaluées en continu :

\begin{table}[H]
\centering
\small
\begin{tabularx}{\linewidth}{@{}l X l@{}}
\toprule
\textbf{Alerte} & \textbf{Condition} & \textbf{Gravité} \\
\midrule
\texttt{APIInjoignable} & Prometheus ne peut plus collecter l'API pendant 1~min & critique \\
\texttt{TauxErreurs5xxEleve} & plus de 5~\% d'erreurs serveur sur 5~min & haute \\
\texttt{LatenceAPIElevee} & p95 de \texttt{/predict} au-delà de 2~s pendant 5~min & moyenne \\
\texttt{RefusAuthentificationNombreux} & plus de 50 refus 401 en 5~min & moyenne \\
\bottomrule
\end{tabularx}
\caption{Règles d'alerte Prometheus de l'API (\texttt{deployment/prometheus/rules/api.yml}).}
\label{tab:alertes-aws}
\end{table}

\paragraph{Exploitation.} La VM s'arrête automatiquement chaque soir à 22~h
(EventBridge Scheduler) et la pile redémarre seule au démarrage suivant. Le disque
a été porté à 80~Go et 4~Go de swap protègent les entraînements d'un manque de
mémoire. Le contexte de construction des images exclut les données et le cache DVC
(\texttt{.dockerignore}) : sans cette exclusion, Docker copiait 12~Go d'images dans
l'image d'entraînement et saturait le disque.

\subsection{Validation de l'API déployée}

Les tests suivants ont été menés sur la VM le 22 octobre 2026, avant le premier
entraînement, puis le 24 octobre après un entraînement de validation du pipeline
(modèle \emph{baseline}, 1\,000 images, 5 époques, F1 macro de test 0,44) promu
\emph{champion} : la chaîne complète, de l'entraînement à la prédiction authentifiée,
fonctionne.

\begin{table}[H]
\centering
\small
\begin{tabularx}{\linewidth}{@{}X c X@{}}
\toprule
\textbf{Requête} & \textbf{Code HTTP} & \textbf{Réponse} \\
\midrule
\texttt{GET /healthz} & 200 & service démarré \\
\texttt{POST /login}, mauvais identifiants & 401 & \texttt{Invalid credentials} \\
\texttt{POST /login}, bons identifiants & 200 & jeton JWT \\
\texttt{POST /predict} sans jeton & 401 & \texttt{Missing authentication token} \\
\texttt{POST /predict}, jeton invalide & 401 & \texttt{Invalid token} \\
\texttt{POST /predict}, jeton valide & 400 & modèle ou alias \emph{champion} absent (attendu) \\
\texttt{GET /models} & 200 & liste vide (attendu) \\
\texttt{GET /metrics} & 200 & compteurs \texttt{blood\_cells\_server\_*} \\
\midrule
\multicolumn{3}{@{}l}{\textit{Après promotion d'un modèle \emph{champion} (24 octobre)}} \\
\texttt{POST /predict}, jeton valide & 200 & classe prédite et probabilités des 10 classes \\
\texttt{GET /models} & 200 & \texttt{baseline-type} \\
\bottomrule
\end{tabularx}
\caption{Tests de l'API sur la VM : l'authentification protège la prédiction, et un modèle promu \emph{champion} est servi.}
\label{tab:tests-api-aws}
\end{table}
